\documentclass[10pt,a4paper]{amsart}
\usepackage[margin=19mm]{geometry}
\usepackage{amsmath,amssymb,mathtools}
\usepackage[scr=rsfs,cal=euler]{mathalfa}
\usepackage{xfrac}
\usepackage[unicode,hidelinks,bookmarks=false]{hyperref}
\newcommand{\ie}{i.e.\ }
\newcommand{\cf}{cf.\ }
\newcommand{\resp}{respectively\ }
\newcommand{\Subsection}{\subsection}
\providecommand{\nobreakdash}{\nobreak\hbox{-}}
\setlength{\emergencystretch}{2em}
\multlinegap0pt
\usepackage{amssymb}
\usepackage{algorithm}\usepackage{algpseudocode}
\usepackage{bm}
\usepackage{dsfont}
\usepackage{enumerate}
\usepackage{tikz}
\usepackage[normalem]{ulem}
\usepackage{xcolor}
\newtheorem{thm}{Theorem}
\title{A Note On Gaussian Random Fields \& Underlying Markov Processes Through a Central Limit Theorem}
\author{Edward C Waymire}
\date{Source: arXiv:2606.21759v3 (CC BY 4.0)}
\begin{document}
\maketitle

\noindent\emph{Source and licence.} A Note On Gaussian Random Fields \& Underlying Markov Processes Through a Central Limit Theorem. Version arXiv:2606.21759v3 (CC BY 4.0), distributed by arXiv under the Creative Commons Attribution 4.0 International licence, http://creativecommons.org/licenses/by/4.0/, as recorded in the arXiv licence metadata for this exact version and shown on the abstract page https://arxiv.org/abs/2606.21759v3. Full licence notice: \texttt{licenses/arxiv-2606.21759-CC-BY-4.0.txt}.\par\smallskip

\noindent\textit{Editorial scope.} Counted unit: the article's thm extendtime (source0.tex lines 792--816). The article's complete original proof of it is delivered with it (source0.tex lines 817--880, one \texttt{proof} environment of 64 lines). No mathematics is written, repaired or re-derived: the statement and the proof are the article's own lines, in the article's own notation, and no retained line is substituted or re-flowed.  No further source-local context is needed: every cross-reference of every retained line resolves inside the two delivered spans.  Nothing was omitted from any retained span, so the package carries no omission record.  No retained line contains a citation, so the wrapper carries no bibliography and no bibliography entry.  No retained line contains a figure environment, an image-inclusion command or drawing code, so no image is delivered and the package declares no asset dependency.  Editorial formatting only, in addition: an AMS \texttt{amsart} preamble that restates the article's own declarations and macros used by the retained lines, namely the environments \texttt{thm} on separate counters, together with the standard packages those lines need.  The retained environments print in this document as Theorem 1 in order of appearance, whereas the article prints the same items with the numbering of its own full document.  The complete native source of the article as distributed in the arXiv e-print for this exact version is bundled unchanged as \texttt{sources/203294/source0.tex}.
\begin{thm}[Brownian Time Extension]\label{extendtime}
Let $\{B(f,t):f\in 1_\pi^\perp, t\in\mathbb{R}\}$ denote the
temporal extension of $B$ with finite dimensional
distributions defined by Gaussian weak-convergence limit 
\begin{equation}
\lim_{n\to\infty}\frac{1}{\sqrt{n}}\bigg(\int_{nt_1}^0f_1(\tilde{X}_{s_1})ds_1,\dots,\int_{nt_k}^0f_k(\tilde{X}_{s_k})ds_k,
 \int_{0}^{nt_{k+1}}f_{k+1}(\tilde{X}_{s_{k+1}})ds_{k+1},\dots,\int_0^{nt_m}f_m(\tilde{X}_{s_m})ds_m\bigg),
\end{equation}
for $t_1<t_2\cdots<t_k<0<t_{k+1}<\cdots< t_m$.
 Then ${B}$ is a space-time
centered 
Gaussian random field indexed by $1_\pi^\perp\times\mathbb{R}$
with 
 \begin{equation}
 \text{cov}(B(f_1,t_1),B(f_2,t_2))
 =
 \begin{cases} (|t_1|\wedge|t_2|)(-2\langle f_1, A^{-1}f_2\rangle_\pi)
 & \text{if}\quad t_1t_2 >0\\
 0 & \text{if}\quad t_1t_2\le 0.
 \end{cases}
 \end{equation}
 Moreover, 
 the path $t\to B(f,t)$ has an
 almost surely continuous version as  Brownian motion.
\end{thm}

\begin{proof}
To simplify the notation involved in the calculations,
let 
\begin{equation}I_a^b(j) = \int_a^bf_j(X_{s_j})ds_j,\quad 
\Delta_a^b(j) = B(f_j,b)-B(f_j,a).\end{equation}
Now, suppose $t_1<t_2\cdots<t_k<0<t_{k+1}<\cdots < t_m$.
For the double-sided extension of $X$, one makes  (integral) changes
 of variable of the form $u_j = s_j-nt_1$ to obtain 
 \begin{eqnarray}
 &&\frac{1}{\sqrt{n}}\bigg(I_{nt_1}^0(1),\dots,I_{nt_k}^0(k),
 I_{0}^{nt_{k+1}}(k+1),\dots,I_0^{nt_m}(m)\bigg)\nonumber\\
 &=&\frac{1}{\sqrt{n}}\bigg(\int_0^{-nt_1}(f_1(\tilde{X}_{u_1+nt_1})du_1,\dots,\int_{n(t_k-t_1)}^{-nt_1}f_k(\tilde{X}_{u_k+nt_1})du_k,\nonumber\nonumber\\
 &&\int_{-{nt_1}}^{n(t_{k+1}-t_1)}f_{k+1}(\tilde{X}_{u_{k+1}+nt_1})du_{k+1},\dots,\int_{-nt_1}^{n(t_m-t_1)}f_m(\tilde{X}_{u_m+nt_1})du_m\bigg)\nonumber\\
&=&\frac{1}{\sqrt{n}}\bigg(I_0^{-nt_1}(1),\dots,I_{n(t_k-t_1)}^{-nt_1}(k),
 I_{-{nt_1}}^{n(t_{k+1}-t_1)}(k+1),\dots,I_{-nt_1}^{n(t_m-t_1)}(m)\bigg)\nonumber\\
&=&\frac{1}{\sqrt{n}}\bigg(I_0^{-nt_1}(1),\dots,
I_0^{-nt_1}(k)
-I_0^{n(t_k-t_1)}(k),
 I_{0}^{n(t_{k+1}-t_1)}(k+1)\nonumber\\
&&-I_0^{-nt_1}(k+1),\dots,
I_{0}^{n(t_m-t_1)}(m)
-I_{0}^{-nt_1}(m)\bigg)\nonumber\\
&&\Rightarrow \bigg(\Delta_0^{-t_1}(1),\dots,
\Delta^{-t_1}_{t_k-t_1}(k),
\Delta^{t_{k+1}-t_1}_{-t_1}(k+1),\dots,
\Delta^{t_m-t_1}_{-t_1}(m)\bigg)\nonumber\\
&=&\bigg(B(f_1,t_1),B(f_2,t_2),\dots,B(f_m,t_m)\bigg). 
 \end{eqnarray}
 That is, the limiting joint finite dimensional
  distribution of $(B(f_1,t_1),\dots,B(f_m,t_m))$ is the Gaussian distribution of 
  $$(\Delta_0^{-t_1}(1),\Delta^{-t_1}_{t_2-t_1}(2),\dots,
  \Delta^{t_m-t_1}_{-t_1}(m))$$
  previously defined for $-t_1\ge 0, t_j-t_1\ge 0, j=2,\dots,m$. 
Thus, computing the expected values
 $\mathbb{E}\Delta^{-t_1}_0(1)\Delta^{-t_1}_{t_j-t_1}(j)$, for $j\le k$, 
 and  $\mathbb{E}\Delta^{-t_1}_0\Delta^{t_j-t_1}_{-t_1}(j)$,
 for  $j\ge k+1$,
 one arrives at 
  \begin{equation}\label{covform1}
\text{cov}(B(f_1,t_1),B(f_j,t_j))
  =\begin{cases} |t_j|
  (-2\langle f_1, A^{-1}f_j\rangle_\pi)\quad j\le k\\
 0\quad k+1\le j\le m.
   \end{cases}
  \end{equation}
Similarly, for $i\ge 2$, computing
 $\mathbb{E}\Delta^{-t_1)}_{t_i-t_1}(i)\Delta^{-t_1}_{t_j-t_1}(j)$, 
 for $2\le i<j\le k$,
 $\mathbb{E}\Delta^{-t_1}_{t_i-t_1}(i)\Delta^{t_j-t_1}_{-t_1}(j)$, 
 for $2\le i\le k<j$, and
 $\mathbb{E}\Delta^{t_i-t_1}_{-t_1)}(i)\Delta^{t_j-t_1}_{-t_1}(j)$, 
 for $k<i<j\le m$,
 one obtains
\begin{equation}\label{covform2}
  \text{cov}(B(f_i,t_i),B(f_j,t_j))
  =
 \begin{cases}|t_j|
  (-2\langle f_1, A^{-1}f_j\rangle_\pi),\  2\le i<j\le k\\
0,\  2\le i <k+1\le j\le m\\
 (t_i\wedge t_j)
  (-2\langle f_1, A^{-1}f_j\rangle_\pi),\ k < i\le j\le m.
  \end{cases}
    \end{equation}
 \end{proof}

\end{document}
